\Needspace{7\baselineskip}
\pdfbookmark[0]{Notices. Müller: Theory of Parallels. 1822.}{anzeigen-mueller-1822}
\begin{center}
\rule{.85\textwidth}{.3pt}\par
Göttingische gelehrte Anzeigen.\quad 28 October 1822.\par
\rule{.85\textwidth}{.3pt}
\end{center}
\hspace*{\parindent}\textit{Theory of Parallels, by} \textsc{Carl Reinhard Müller}, \textit{Doctor of Philosophy, extraordinary professor of mathematics, and so on.} 40 pages in quarto. Marburg 1822.

Six years ago the reviewer already expressed his conviction in these pages that all previous attempts to prove the theory of parallel lines rigorously or to fill the gap in \textsc{Euclid}ean geometry had brought us no closer to that goal. He can only extend this judgment to all subsequent attempts that have come to his attention. Even so, despite missing their real aim, some of these attempts are still worth reading for anyone fond of geometry because of the ingenuity they show, and. In this respect the reviewer believes he must single out the present short work, published on the occasion of a school examination, for particular recognition. To set out the author's whole ingenious line of thought in detail here would take up too much room in these pages and would also be unnecessary, since the work itself deserves to be read. It has a weak point, however, as all the other attempts do, and the purpose of this review is to bring it out. We find this weak point on p.~15 in the proof of the theorem of Article~15. This theorem is the vital nerve of the entire theory, which collapses if the theorem cannot be proved rigorously. We therefore begin by stating the theorem here. Anyone can easily draw the corresponding figure.

If each angle at the base \(ON\) of an isosceles triangle is greater than the angle at the vertex \(A\), and if at \(O\) one constructs on the side \(OA\) an angle equal to angle \(A\) whose other arm \(OL\) meets \(AN\) at the point \(L\) between \(A\) and \(N\), and then cuts off
